\documentclass[12pt]{article}
\usepackage[utf8]{inputenc}
\usepackage[T1]{fontenc}
\usepackage[margin=1in]{geometry}
\usepackage{setspace}
\usepackage{xcolor}
\usepackage[colorlinks=true, urlcolor=blue!60!black, linkcolor=blue!60!black]{hyperref}
\setlength{\parindent}{2em}
\setlength{\parskip}{1.25em}
\renewcommand{\baselinestretch}{1.0}
\begin{document}
\begin{center}
{\LARGE \textit{Baker v.\ Carr}, 369 U.S.\ 186 (1962)\par}
\vspace{0.8em}
{\large Case Brief --- Four-Page Edition\par}
\end{center}
\vspace{1em}

\section{Memory Jogger}

Equal Protection challenges to malapportioned state legislative districts are justiciable in federal court; the Court also articulated the modern six-factor test governing the political question doctrine (\href{https://supreme.justia.com/cases/federal/us/369/186/}{Justia}).

\section{Detailed Case Facts}

Tennessee's constitution required the General Assembly to reapportion its 33 Senate and 99 House seats every ten years according to the number of qualified voters. The legislature complied through 1901 but enacted no general reapportionment after the Apportionment Act of 1901; every later reapportionment bill failed. Meanwhile the population grew and urbanized, and the frozen map produced severe disparities: about 37\% of voters elected 20 of 33 Senators, and 40\% elected 63 of 99 House members.

Charles W. Baker and other urban voters sued Tennessee officials under 42 U.S.C. \S~1983, alleging a ``debasement of their votes'' in violation of the Equal Protection Clause. Tennessee offered no initiative or referendum, and the Tennessee Supreme Court had declined relief in \textit{Kidd v. McCanless} (1956), leaving urban voters with no effective state remedy. Defendants argued that the claim was a political question and sounded only under the nonjusticiable Guaranty Clause.

\section{Procedural History}

A three-judge district court dismissed for lack of subject-matter jurisdiction and failure to state a claim. It agreed that the legislature was ``guilty of a clear violation of the state constitution and of the rights of the plaintiffs,'' but held that \textit{Colegrove v. Green}, 328 U.S. 549 (1946), foreclosed federal intervention. On direct appeal, after argument and reargument (with Solicitor General Archibald Cox as amicus supporting reversal), the Supreme Court reversed and remanded on March 26, 1962.

\section{Judicial Votes}

Reversed and remanded, \textbf{6--2}.
\begin{itemize}
\item \textbf{Majority:} Justice Brennan, joined by Chief Justice Warren and Justices Black, Douglas, Clark, and Stewart.
\item \textbf{Concurrences:} Justices Douglas, Clark, and Stewart, each joining the opinion of the Court.
\item \textbf{Dissents:} Justice Frankfurter, joined by Justice Harlan; Justice Harlan separately.
\item \textbf{Not participating:} Justice Whittaker.
\end{itemize}

\section{Holding}

(1) Federal courts have subject-matter jurisdiction over a non-frivolous claim that a state apportionment statute violates the Fourteenth Amendment; (2) voters alleging arbitrary vote dilution have standing; and (3) an Equal Protection challenge to state legislative malapportionment is justiciable, not a political question. A case should not be dismissed as a political question unless one of six attributes is ``inextricable from the case at bar.'' The Court did not decide the merits; the substantive standard (``one person, one vote'') came later in \textit{Wesberry v. Sanders} and \textit{Reynolds v. Sims} (1964).

\section{Analysis of Opinions}

\textbf{Majority (Brennan).} Brennan separated jurisdiction, standing, and justiciability, which the district court had conflated. Jurisdiction existed because the claim was not frivolous; plaintiffs had standing because they asserted an individual injury, the impairment of their votes. The political question doctrine, he reasoned, concerns ``the relationship between the judiciary and the coordinate branches of the Federal Government,'' not the federal judiciary's relationship to the States. Surveying foreign affairs, war, enactment validity, Indian tribes, and the Guaranty Clause, he distilled six factors. Guaranty Clause claims under \textit{Luther v. Borden} are nonjusticiable only because they implicate those factors; this Equal Protection claim involved no coordinate federal branch and invoked ``well developed and familiar'' standards, as in \textit{Gomillion v. Lightfoot} (1960).

\textbf{Douglas (concurring).} Voting rights are foundational, and a 19-to-1 disparity, if proved, would establish a clear violation.

\textbf{Clark (concurring).} The Court should have reached the merits: Tennessee's apportionment was ``a crazy quilt without rational basis.'' Because voters were ``caught up in a legislative strait jacket,'' courts were the only avenue of relief. \textbf{Stewart (concurring)} stressed that the Court did not require ``approximate equality'' of every vote.

\textbf{Frankfurter (dissenting).} The decision was ``a massive repudiation of the experience of our whole past.'' The claim was a ``Guarantee Clause claim masquerading under a different label,'' requiring a choice among theories of representation with no judicial standards; relief belonged to ``an informed, civically militant electorate.''

\textbf{Harlan (dissenting).} Nothing in the Fourteenth Amendment requires legislatures ``to reflect with approximate equality the voice of every voter''; weighting geography is a legislative judgment.

\section{Examples: Future Applications}

\begin{enumerate}
\item \textbf{Same-Side (Justiciable): Passport Statute.} A statute directs the State Department to record ``Israel'' on passports of citizens born in Jerusalem. Under \textit{Zivotofsky v. Clinton} (2012), deciding the statute's constitutionality is a familiar judicial task, not foreign policymaking.
\item \textbf{Same-Side (Justiciable): County-Based Malapportionment.} Alabama guarantees each county a representative, producing a 41-to-1 disparity. Justiciable and, under \textit{Reynolds v. Sims}, unconstitutional.
\item \textbf{Opposite-Side (Nonjusticiable): Senate Impeachment Procedure.} A judge challenges the Senate's use of a trial committee. Under \textit{Nixon v. United States} (1993), ``sole Power to try'' is a textual commitment and ``try'' lacks manageable standards.
\item \textbf{Opposite-Side (Nonjusticiable): Partisan Gerrymandering.} In \textit{Rucho v. Common Cause} (2019), the Court held that how much partisanship is ``too much'' lacks a judicially manageable standard.
\item \textbf{Fence-Sitter: War Powers Suit.} Members of Congress sue to end airstrikes lasting past 60 days. Construing a statute is judicial work, but whether ``hostilities'' exist implicates Factors 1, 2, and 6, and lower courts often dismiss such suits.
\end{enumerate}

\section{Critique}

\textbf{An under-theorized test.} Brennan purported to summarize past cases but transformed the doctrine into a separation-of-powers-only inquiry, while his claim that ``well developed and familiar'' standards existed was largely ipse dixit; \textit{Reynolds} soon had to adopt a bright-line rule that Stewart had disclaimed. Originalists, echoing Harlan, note that Section 2 of the Fourteenth Amendment contemplated state structuring of the franchise, and question the prudential factors 4 through 6.

\textbf{Representation reinforcement.} John Hart Ely and other process theorists celebrate \textit{Baker}: when insiders lock the channels of change, abstention entrenches minority rule, and Frankfurter's appeal to the electorate was circular. Yet \textit{Baker}'s ``judicially manageable standards'' prong later insulated partisan gerrymandering in \textit{Rucho}.

\section{Key Quotations}

\begin{quote}
``Prominent on the surface of any case held to involve a political question is found a textually demonstrable constitutional commitment of the issue to a coordinate political department; or a lack of judicially discoverable and manageable standards for resolving it; or the impossibility of deciding without an initial policy determination of a kind clearly for non judicial discretion; or the impossibility of a court's undertaking independent resolution without expressing lack of the respect due coordinate branches of government; or an unusual need for unquestioning adherence to a political decision already made; or the potentiality of embarrassment from multifarious pronouncements by various departments on one question.''
\end{quote}
\noindent --- Justice Brennan, for the Court, 369 U.S. at 217.

\begin{quote}
``Have the appellants alleged such a personal stake in the outcome of the controversy as to assure that concrete adverseness which sharpens the presentation of issues upon which the court so largely depends for illumination of difficult constitutional questions? This is the gist of the question of standing.''
\end{quote}
\noindent --- Justice Brennan, for the Court, 369 U.S. at 204.

\begin{quote}
``The doctrine of which we treat is one of `political questions,' not one of `political cases.' The courts cannot reject as `no law suit' a bona fide controversy as to whether some action denominated `political' exceeds constitutional authority.''
\end{quote}
\noindent --- Justice Brennan, for the Court, 369 U.S. at 217.

\end{document}
